\sectioncentered*{Перечень условных обозначений}
\addcontentsline{toc}{section}{Перечень условных обозначений}
\label{sec:reduction}
\begin{flushleft}
	\begin{tabular}{@{}l@{\hspace{0.4em}}c@{\hspace{0.4em}}p{11.2cm}@{}}
		TF & \mdash & Частота термина в документе или предложении \\
		IDF & \mdash & Обратная документная частота \\
		TF-IDF & \mdash & Вес термина по частоте и редкости в коллекции \\
		Posd & \mdash & Функция положения предложения в документе \\
		Posp & \mdash & Функция положения предложения в абзаце \\
		Score & \mdash & Суммарный вес слов предложения \\
		OSTIS & \mdash & Open Semantic Technology for Intelligent Systems \\
		F-мера & \mdash & Гармоническое среднее точности и полноты \\
		PDF & \mdash & Переносимый формат документов \\
		API & \mdash & Программный интерфейс приложения \\
		REST & \mdash & Архитектурный стиль взаимодействия компонентов \\
		JSON & \mdash & Текстовый формат обмена данными \\
		SQL & \mdash & Язык структурированных запросов \\
		СУБД & \mdash & Система управления базами данных \\
		SPA & \mdash & Одностраничное веб-приложение \\
	\end{tabular}
\end{flushleft}
